% !TeX program = xelatex
% NJU-LINK technical report. Recommended Overleaf compiler: XeLaTeX.
% The original FAIR / TokenWave front-matter commands are preserved.
% Change the metadata below and replace the demonstration text.
\documentclass[10pt]{njulink}

\njulinksetup{
  type={Technical Report},
  id={NJU-LINK / RESEARCH SERIES},
  short-title={NJU-LINK Technical Report},
  subtitle={A clear format for ideas, methods, and evidence},
  keywords={Research reports; reproducibility; scientific communication}
}

\title{NJU-LINK Technical Report}
\author[1]{First Author}
\author[1,2]{Second Author}
\author[1]{Third Author}
\affiliation[1]{NJU-LINK}
\affiliation[2]{Affiliation Two}
% \contribution[*]{Equal contribution}
% \correspondence{First Author at \email{name@example.org}}
% \metadata[Code]{\url{https://example.org/project}}
\date{October 2026}

\abstract{%
A research report should make its question, method, and evidence easy to find.
This template provides a consistent structure for NJU-LINK technical reports,
working papers, and research notes. A compact identity block, a restrained
blue--violet palette, and a readable text column establish the visual hierarchy.
The following pages demonstrate headings, author affiliations, citations,
a vector figure, an equation, and a comparison table. This is a typesetting
example; it does not describe an empirical study or report experimental results.
}

\begin{document}
\maketitle

\section{Introduction}
\label{sec:introduction}
Start with the research problem and the observation that makes it worth
investigating. Define the setting precisely enough for the reader to understand
what is being measured, which assumptions matter, and where the proposed
approach is expected to apply. A concrete example often communicates the
motivation more effectively than a broad statement about the field.

Then connect the limitation of current practice, the idea that addresses it,
and the evidence used to evaluate that idea. State each contribution at the
level supported by the results. The method appears in \Cref{sec:method};
\Cref{tab:reporting} summarizes what a comparison should specify.

\begin{njubox}[title={The central question}]
What changes, why should that change help, and which observation would show
that it works? Use this space for one research question or one supported
takeaway, then develop the evidence in the body of the report.
\end{njubox}

\section{Related Work}
Organize prior work by the distinction relevant to the research question.
Explain which assumptions, capabilities, or evaluation settings differ.
For example, \citet{vaswani2017attention} can be cited as part of a sentence;
a parenthetical citation uses the complementary form
\citep{vaswani2017attention}. This reference demonstrates the bibliography style.

\section{Method}
\label{sec:method}
Define the inputs, outputs, and assumptions, then describe the stages
in \Cref{fig:workflow} in the order needed to reproduce the method.

\begin{figure}[htbp]
\centering
\begin{tikzpicture}[
  stage/.style={draw=LinkRule,line width=0.6pt,rounded corners=6pt,
    minimum width=34mm,minimum height=17mm,align=left,inner sep=8pt,
    text width=27mm,font=\sffamily\small,text=LinkNavy},
  connector/.style={-{Stealth[length=2mm,width=1.4mm]},
    draw=LinkMuted,line width=0.7pt}
]
  \node[stage,fill=LinkPanel] (question) at (0,0)
    {{\color{LinkBlue}\scriptsize\bfseries 01}\quad\textbf{Question}\\[4pt]
     {\scriptsize Scope and assumptions}};
  \node[stage,fill=LinkLavender] (method) at (4.05,0)
    {{\color{LinkPurple}\scriptsize\bfseries 02}\quad\textbf{Method}\\[4pt]
     {\scriptsize Inputs and mechanism}};
  \node[stage,fill=LinkMint] (evidence) at (8.10,0)
    {{\color{LinkTeal}\scriptsize\bfseries 03}\quad\textbf{Evidence}\\[4pt]
     {\scriptsize Controls and analysis}};
  \node[stage,fill=LinkSand] (release) at (12.15,0)
    {{\color{LinkMuted}\scriptsize\bfseries 04}\quad\textbf{Release}\\[4pt]
     {\scriptsize Code and limitations}};
  \draw[connector] (question.east) -- (method.west);
  \draw[connector] (method.east) -- (evidence.west);
  \draw[connector] (evidence.east) -- (release.west);
\end{tikzpicture}
\caption{An illustrative research workflow. Colors identify stages and do not encode measured values.}
\label{fig:workflow}
\end{figure}

\subsection{Notation and evaluation}
For $N$ examples with individual scores $s_i$, the arithmetic mean is
\begin{equation}
  \overline{s}=\frac{1}{N}\sum_{i=1}^{N}s_i.
  \label{eq:mean}
\end{equation}
For \Cref{eq:mean}, specify the score's unit and the handling of missing observations.

\section{Experiments}
State the data, baselines, and controlled variables before presenting results.
\Cref{tab:reporting} demonstrates the table style without synthetic performance numbers.

\begin{table}[htbp]
\caption{A compact reporting checklist. Replace these rows with the study's own comparisons.}
\label{tab:reporting}
\centering\small
\begin{tabularx}{\linewidth}{lYY}
\toprule
\rowcolor{LinkPanel}
\textbf{Element} & \textbf{What to specify} & \textbf{Why it matters} \\
\midrule
Data & Source, split, and exclusions & Defines the evaluated population \\
Comparison & Baselines and shared settings & Makes differences interpretable \\
Measurement & Metric, unit, and aggregation & Connects values to the question \\
Uncertainty & Repeats and variation & Shows the stability of the result \\
\bottomrule
\end{tabularx}
\end{table}

\section{Results and Analysis}
Pair each claim with the comparison that supports it. Distinguish a measured
difference from a proposed explanation of that difference.

\section{Limitations}
Describe the conditions under which the finding holds, the settings not
evaluated, and the comparisons still needed.

\section{Conclusion}
Close with the supported finding, its practical implication, and the next experiment.

% Add acknowledgments only when there is something to acknowledge.
% \section*{Acknowledgments}

\bibliographystyle{assets/plainnat}
\bibliography{references}

\clearpage
\beginappendix
\section{Reproducibility Details}
\label{app:details}
Use the appendix for configurations, additional analyses, and details needed
to reproduce the study. Appendix sections use lettered numbering; they can be
referenced from the main text with the same cross-reference commands.

\subsection{Configuration record}
Record the information required to connect the result to a specific run:
\begin{enumerate}
  \item the version of the data and the rule used to construct each split;
  \item the implementation version, configuration, and random seeds;
  \item the evaluation script, metric definitions, and aggregation procedure.
\end{enumerate}

\begin{njubox}[title={Using this template}]
Edit the title, authors, affiliations, and abstract at the beginning of
\texttt{main.tex}. The original author, affiliation, abstract, and metadata
commands are retained. The short Chinese starting point is \texttt{starter.tex};
the README explains the available settings.
\end{njubox}

\subsection{Supporting material}
Provide enough detail for another researcher to inspect the assumptions and
repeat the analysis. Prefer descriptive figure captions and explicit units;
use tables for comparisons and short paragraphs for interpretation.
\end{document}
